\section{Bases}
\label{section-bases}

\noindent
Basic material on bases for topological spaces.

\begin{definition}
\label{definition-base}
Let $X$ be a topological space. A collection of subsets $\mathcal{B}$ of $X$
is called a {\it base for the topology on $X$} or a {\it basis for the
topology on $X$} if the following conditions hold:
\begin{enumerate}
\item Every element $B \in \mathcal{B}$ is open in $X$.
\item For every open $U \subset X$ and every $x \in U$,
there exists an element $B \in \mathcal{B}$ such that
$x \in B \subset U$.
\end{enumerate}
\end{definition}

\noindent
The following lemma is sometimes used to define a topology.

\begin{lemma}
\label{lemma-make-base}
Let $X$ be a set and let $\mathcal{B}$ be a collection of subsets.
Assume that $X = \bigcup_{B \in \mathcal{B}} B$ and that given
$x \in B_1 \cap B_2$ with $B_1, B_2 \in \mathcal{B}$ there is a
$B_3 \in \mathcal{B}$ with $x \in B_3 \subset B_1 \cap B_2$.
Then there is a unique topology on $X$ such that $\mathcal{B}$
is a basis for this topology.
\end{lemma}

\begin{proof}
Let $\sigma(\mathcal B)$ be the set of subsets of $X$ which can be written
as unions of elements of $\mathcal B$. We claim $\sigma(\mathcal{B})$
is a topology. Namely, the empty set is an element of $\sigma(\mathcal{B})$
(as an empty union) and $X$ is an element of $\sigma(\mathcal{B})$ (as the
union of all elements of $\mathcal{B}$). It is clear that $\sigma(\mathcal{B})$
is preserved under unions. Finally, if $U, V \in \sigma(\mathcal{B})$
then write $U = \bigcup_{i \in I} U_i$ and $V = \bigcup_{j \in J} V_j$
with $U_i, V_j \in \mathcal{B}$ for all $i \in I$ and $j \in J$.
Then
$$
U \cap V = \bigcup\nolimits_{i \in I, j \in J}
U_i \cap V_j
$$
The assumption in the lemma tells us that $U_i \cap V_j \in \sigma(\mathcal{B})$
hence we see that $U \cap V$ is too. Thus $\sigma(\mathcal{B})$ is a topology.
Properties (1) and (2) of Definition \ref{definition-base} are
immediate for this topology. To prove the uniqueness of this topology
let $\mathcal T$ be a topology on $X$ such that $\mathcal B$ is a base for it.
Then of course every element of $\mathcal{B}$ is in $\mathcal{T}$ by
(1) of Definition \ref{definition-base} and
hence $\sigma(\mathcal{B}) \subset \mathcal{T}$. Conversely, part
(2) of Definition \ref{definition-base} tells us that every element
of $\mathcal{T}$ is a union of elements of $\mathcal{B}$, i.e.,
$\mathcal{T} \subset \sigma(\mathcal{B})$. This finishes the proof.
\end{proof}

\begin{lemma}
\label{lemma-refine-covering-basis}
Let $X$ be a topological space.
Let $\mathcal{B}$ be a basis for the topology on $X$.
Let $\mathcal{U} : U = \bigcup_i U_i$ be an open covering of
$U \subset X$. There exists an open covering $U = \bigcup V_j$
which is a refinement of $\mathcal{U}$ such that each
$V_j$ is an element of the basis $\mathcal{B}$.
\end{lemma}

\begin{proof}
If $ x \in U = \bigcup_{i\in I} U_i $, there is an $ i_x \in I $ such that
$ x \in U_{i_x} $. Thus we have a $ B_{i_x} \in \mathcal{B}$
verifying $ x \in B_{i_x} \subset U_{i_x}$. Set
$J = \{i_x | x \in U\}$ and for $j = i_x \in J$ set $V_j = B_{i_x}$.
This gives the desired open covering of $U$ by $\{V_j\}_{j \in J}$.
\end{proof}

\begin{definition}
\label{definition-subbase}
Let $X$ be a topological space. A collection of subsets $\mathcal{B}$ of $X$
is called a {\it subbase for the topology on $X$} or a {\it subbasis for the
topology on $X$} if the finite intersections of
elements of $\mathcal{B}$ form a basis for the topology on $X$.
\end{definition}

\noindent
In particular every element of $\mathcal{B}$ is open.

\begin{lemma}
\label{lemma-subbase}
Let $X$ be a set. Given any collection $\mathcal{B}$ of subsets of $X$
there is a unique topology on $X$ such that $\mathcal{B}$ is a subbase
for this topology.
\end{lemma}

\begin{proof}
By convention $\bigcap_\emptyset B = X $. Thus we can apply Lemma
\ref{lemma-make-base} to the set of finite intersections of elements
from $\mathcal B$.
\end{proof}

\begin{lemma}
\label{lemma-create-map-from-subcollection}
Let $X$ be a topological space. Let $\mathcal{B}$ be a collection
of opens of $X$. Assume $X = \bigcup_{U \in \mathcal{B}} U$ and
for $U, V \in \mathcal{B}$ we have
$U \cap V = \bigcup_{W \in \mathcal{B}, W \subset U \cap V} W$.
Then there is a continuous map $f : X \to Y$ of topological spaces
such that
\begin{enumerate}
\item for $U \in \mathcal{B}$ the image $f(U)$ is open,
\item for $U \in \mathcal{B}$ we have $f^{-1}(f(U)) = U$, and
\item the opens $f(U)$, $U \in \mathcal{B}$
form a basis for the topology on $Y$.
\end{enumerate}
\end{lemma}

\begin{proof}
Define an equivalence relation $\sim$ on points of $X$
by the rule
$$
x \sim y \Leftrightarrow
(\forall U \in \mathcal{B} : x \in U \Leftrightarrow y \in U)
$$
Let $Y$ be the set of equivalence classes and $f : X \to Y$
the natural map. Part (2) holds by construction.
The assumptions on $\mathcal{B}$ exactly
mirror the assumptions in Lemma \ref{lemma-make-base}
on the set of subsets $f(U)$, $U \in \mathcal{B}$.
Hence there is a unique topology on $Y$ such that (3) holds.
Then (1) is clear as well.
\end{proof}
